\scanpage{78}
\mspage{\handwritten{- 1 -}}

\noindent\uline{A\,V\,A\,N\,T\,-\,P\,R\,O\,P\,O\,S}

\medskip
\noindent\begin{minipage}{\linewidth}\small
I.\qquad
\tikz[remember picture,overlay]{\coordinate (cotesixtythreestart) at (-0.7em,1.8ex);}
Le but principal du présent Séminaire est de développer le formalisme de la ``cohomologie de Weil'' des Schémas. A partir essentiellement des résultats qui sont démontrés ici, des arguments bien connus, d'ailleurs dûs à Weil lui-même, \'{p}ermettent de déduire une partie des conjectures de Weil sur les fonctions $\mathrm{L}$ des variétés projectives non singulières sur un corps fini\textsuperscript{\handwritten{(1)}}. Cette déduction sera d'ailleurs présentée dans SGA \handwritten{1964/65}, dans un séminaire qui fera suite à celui-ci.

Dans le présent séminaire, nous nous bornons à l'étude de la cohomologie des schémas relativement à la \uline{topologie étale} . Cette topologie par sa description est fort proche des topologies des variétés topologiques habituelles\handwritten{,} et on verra que la plupart des résultats classiques concernant la cohomologie des espaces topologiques.ordinaires (suites spectrales variées, théorèmes de finitude, K\handwritten{\"}{u}nneth, dualité, théorème de Lefschetz) peuvent se formuler et se démontrer dans le nouveau contexte, à condition de se borner le cas échéant aux faisceaux de torsion \uline{premiers aux caractéristiques résiduelles} des schémas envisagés\unsure{.|,|;} On obtiendra une théorie cohomologique ``à coefficient de caractéristique $0$'' (comme demandée par Weil) par un passage à la limite projective e\makebox[0pt][l]{\deleted{\unsure{ss|xx}}}ssentiellement trivial (cf SGA \handwritten{1965}), permettant de définir une cohomologie à coefficient d\deleted{\unsure{e}}\handwritten{\unsure{a|e}ns} l'anneau $\underline{\mathrm{Z}_{\mathord{\rlap{\handwritten{$\scriptstyle/$}}\mathrm{l}}}}$ des entiers $\mathord{\rlap{\handwritten{$/$}}\mathrm{l}}$-adiques à partir des coefficients $\underline{\mathrm{Z}}\handwritten{/}\mathord{\rlap{\handwritten{$/$}}\mathrm{l}}^{\handwritten{\nu}}\underline{\mathrm{Z}}$ ., $\handwritten{\nu\longrightarrow}+\handwritten{\infty}$ . Lorsque $\mathord{\rlap{\handwritten{$/$}}\mathrm{l}}$ est premier aux caractéristiques résiduelles, cette cohomologie possède toutes les bonnes propriétés habituelles dans la cohomologie à coefficient $\underline{\mathrm{Z}}$ classique (et se prête donc à la formulation des conjectures de Weil). D'ailleurs, lorsque les schémas envisagés sont de type fini sur le corps des complexes $\underline{\mathrm{C}}$, nous verrons que la cohomologie (à coefficient\handwritten{s} de torsion) est essentiellement identique à la cohomologie

\smallskip
\noindent\rule{0.32\linewidth}{0.4pt}\par
{\footnotesize\noindent\textsuperscript{\handwritten{(1)}}\quad Au moment d'écrire ces lignes, il n'est pas prouvé que les valeurs propres de l'homomorphisme de Frobenius opérant sur les $\mathrm{H}^{\handwritten{i}}\handwritten{(}\mathrm{X},\underline{\mathrm{Z}}_{\handwritten{\ell}}\handwritten{)}$ sont des entiers algébriques, ni à fortiori que leurs\handwritten{\tikz[remember picture,overlay]{\draw[line width=0.4pt] (cotesixtythreestart) -- (-0.45em,0.15ex);}} valeurs absolues sont égales à $\mathrm{q}^{\handwritten{i/2}}$ .\par}
\end{minipage}
% archive sheet number: 77
